\documentclass[conference]{IEEEtran}
\IEEEoverridecommandlockouts
\usepackage{cite}
\usepackage{amsmath,amssymb,amsfonts}
\usepackage{graphicx}
\usepackage{textcomp}
\usepackage{xcolor}
\usepackage{array}
\def\BibTeX{{\rm B\kern-.05em{\sc i\kern-.025em b}\kern-.08em
    T\kern-.1667em\lower.7ex\hbox{E}\kern-.125emX}}
\begin{document}

% ============================================================
%  TITLE
% ============================================================
\title{A Hardware-Enforced Self-Destructing Biometric Template Vault
with Quality-Adaptive Dynamic Threshold Matching in Verilog}

% ============================================================
%  AUTHORS
% ============================================================
\author{
\IEEEauthorblockN{1\textsuperscript{st} Chhavi Talele}
\IEEEauthorblockA{\textit{Dept. of Electronics and Computer Science Engineering} \\
\textit{Vidyalankar Institute of Technology}\\
Mumbai, India \\
chhavi.talele@vit.edu.in}
\and
\IEEEauthorblockN{2\textsuperscript{nd} Malhar Kulkarni}
\IEEEauthorblockA{\textit{Dept. of Electronics and Computer Science Engineering} \\
\textit{Vidyalankar Institute of Technology}\\
Mumbai, India \\
malhar.kulkarni@vit.edu.in}
\and
\IEEEauthorblockN{3\textsuperscript{rd} Hardik Gawande}
\IEEEauthorblockA{\textit{Dept. of Electronics and Computer Science Engineering} \\
\textit{Vidyalankar Institute of Technology}\\
Mumbai, India \\
hardik.gawande@vit.edu.in}
}

\maketitle

% ============================================================
%  ABSTRACT
% ============================================================
\begin{abstract}
Biometric authentication at the hardware level faces two fundamental
problems that software alone cannot solve. First, once a fingerprint
template is stored inside a chip, it stays there permanently---even
after multiple failed attacks---giving an adversary unlimited time to
extract it through physical probing. Second, existing hardware matchers
use a single fixed acceptance threshold, which either incorrectly
rejects genuine users with wet or dry fingers, or incorrectly admits
impostors when the threshold is relaxed. This paper presents a
complete Verilog System-on-Chip (SoC) that solves both problems
directly in hardware, without relying on any external processor or
software. The first contribution is a self-destructing template vault
that automatically wipes all stored fingerprint data using a
three-pass overwrite sequence after three failed login attempts, a
physical tamper signal, or an inactivity timeout---permanently locking
the chip. The second contribution is a quality-adaptive threshold
engine that reads the live scan quality of each fingerprint in real
time and adjusts how strictly it compares the probe against the stored
template. A clean, dry scan is matched against a tight threshold while
a wet or degraded scan is matched against a relaxed threshold,
eliminating unnecessary false rejections while maintaining security.
The design is implemented across six synthesizable Verilog modules,
simulated using Icarus Verilog across seven test scenarios, and
synthesized on the AMD Artix-7 FPGA achieving approximately 175 MHz
operation at 62 mW total power using only 82 look-up tables and
51 flip-flops with zero DSP blocks or block RAM.
\end{abstract}

\begin{IEEEkeywords}
biometric authentication, fingerprint matching, Verilog, FPGA, VLSI,
hardware security, template zeroization, adaptive threshold, quality
analysis, finite state machine, LFSR, anti-tamper, SoC
\end{IEEEkeywords}

% ============================================================
%  SECTION I: INTRODUCTION
% ============================================================
\section{Introduction}

Fingerprint recognition is one of the most widely deployed forms of
biometric authentication, used in smartphones, access control systems,
banking terminals, and IoT devices. When fingerprint matching is
performed in software on a general-purpose processor, the enrolled
template data and the matching algorithm are both accessible to
software-level attackers. Moving the entire authentication pipeline
into dedicated hardware removes this software attack surface entirely
and enables significantly faster, lower-power, and more tamper-resistant
operation~\cite{verilog_biometric}.

However, existing hardware implementations of fingerprint matchers
leave two critical security gaps unaddressed.

The first gap concerns \textit{template permanence}. In every published
FPGA and ASIC biometric implementation, the enrolled fingerprint
template is stored in memory and stays there indefinitely. If an
attacker gains physical access to the chip---through invasive probing,
power-glitching, or decapsulation---the template data can be read out
intact. Unlike a stolen password, a stolen fingerprint cannot be
changed. Prior work on Physical Unclonable Functions
(PUFs)~\cite{puf_auth} has explored binding device identity to
hardware characteristics, and research on secret key extraction from
integrated circuits~\cite{secret_key_ic} has shown that volatile keys
derived from silicon manufacturing variation can be made to
self-extinguish when power is removed. However, no prior work applies
this principle to biometric templates specifically: implementing an
autonomous hardware mechanism that, upon detecting an attack,
immediately destroys all enrolled biometric data using a
standards-compliant multi-pass overwrite.

The second gap concerns \textit{fixed threshold matching}. All
published hardware biometric matchers compare a fingerprint scan
against an enrolled template using a single, unchanging threshold
value. If the threshold is set too strictly, genuine users with
slightly sweaty, dry, or scarred fingers are incorrectly rejected
(False Rejection). If the threshold is set too loosely, the system
admits impostors more easily (False Acceptance). This tradeoff cannot
be resolved with a fixed threshold. Prior Verilog biometric work,
including DES-based biometric co-processors~\cite{verilog_biometric},
reports this constraint without resolving it. The solution requires
the hardware to measure how good each scan is in real time, and adjust
the acceptance decision accordingly.

This paper addresses both gaps with a complete, synthesizable
Verilog SoC. The specific contributions are:

\begin{enumerate}
    \item \textbf{Self-Destructing Template Vault:} A hardware FSM
    that executes a three-pass NIST SP 800-88 compliant memory
    overwrite of all enrolled templates upon detecting three
    consecutive failed authentication attempts, assertion of a
    dedicated tamper input pin, or expiration of an inactivity timer.
    After the wipeout completes, the chip enters a permanent hardware
    lockdown with no recovery path. This is the first published
    hardware implementation of autonomous biometric template
    zeroization.

    \item \textbf{Quality-Adaptive Dynamic Threshold:} A combinational
    hardware engine that fuses three live sensor quality signals---
    surface noise level, ridge contrast, and spatial uniformity---into
    a single quality score, and uses that score to compute the
    matching acceptance threshold in real time. The threshold tightens
    for clean, pristine scans and relaxes for noisy or wet scans,
    eliminating false rejections without compromising security. This is
    the first published hardware implementation of scan quality-driven
    threshold adaptation.

    \item \textbf{Anti-Spoof Guard:} A hardwired minimum contrast
    check that immediately rejects any scan where the ridge contrast
    falls below a set floor value, catching flat artificial fingers
    (gelatin, silicone, latex) before the matcher even runs.

    \item \textbf{Complete Verilog SoC:} A full six-module design
    verified by a self-checking testbench across seven attack and
    environmental scenarios, and synthesized on AMD Artix-7 FPGA.
\end{enumerate}

% ============================================================
%  SECTION II: BACKGROUND AND RELATED WORK
% ============================================================
\section{Background and Related Work}

\subsection{Verilog-Based Biometric Hardware}

A key reference for this work is the Verilog implementation of a
biometric-based secure data transmission system reported
in~\cite{verilog_biometric}. That work implements a Dual-Tree Complex
Wavelet Transform (DT-CWT) for multi-focus image fusion and a
biometric-keyed Data Encryption Standard (DES) co-processor in which
the user's raw fingerprint or iris scan supplies the 56-bit encryption
key. The design uses eight DES substitution boxes (S-boxes) implemented
in Verilog \texttt{case} statements and achieves a total on-chip power
of 0.184~W on Xilinx ISE 14.2. That system demonstrates that biometric
data can be effectively integrated into hardware security pipelines,
but like all prior work, it stores the enrolled biometric template in
permanent memory and applies a fixed matching criterion. Our work
directly extends this direction by adding autonomous template
destruction and adaptive matching thresholds.

\subsection{Physical Unclonable Functions for Authentication}

Physical Unclonable Functions (PUFs) exploit unavoidable variations in
semiconductor manufacturing to generate device-unique digital
signatures without storing a key in memory~\cite{puf_auth}. Arbiter
PUFs route a challenge signal through two symmetric delay paths and
measure which arrives first; Ring Oscillator PUFs compare the
frequencies of identically-laid-out inverter rings to generate
challenge-response bits. Because the PUF response is derived from
physical variation rather than stored data, it cannot be extracted by
reading memory. Our self-destructing vault draws conceptual inspiration
from this idea: rather than using silicon physics to avoid storing
secrets, we implement an active hardware mechanism that destroys stored
secrets when an attack is detected.

\subsection{Secret Key Extraction from Integrated Circuits}

Research on secret key derivation from integrated circuits~\cite{secret_key_ic}
has established that volatile keys tied to silicon characteristics
can be made ephemeral---they exist only when the circuit is powered
and under the correct environmental conditions. This principle
motivates our Time-to-Live (TTL) mechanism: the enrolled template
is treated as an ephemeral credential that expires after a configurable
period of inactivity, after which the hardware autonomously purges it
without any software command.

\subsection{Hardware Security and Fault Attacks}

Research on fault injection attacks against cryptographic
hardware~\cite{fault_attack} demonstrates that adversaries can
deliberately introduce hardware faults during computation to
extract secret values. These attacks are particularly powerful against
systems that retain their secrets indefinitely. Our three-pass
zeroization sequence and permanent lockdown state are specifically
designed as countermeasures: once an anomaly is detected, the secret
is immediately and irreversibly destroyed before any fault-induced
information leakage can occur.

% ============================================================
%  SECTION III: SYSTEM ARCHITECTURE
% ============================================================
\section{Proposed System Architecture}

The proposed system is organized as a six-module Verilog SoC with a
dedicated testbench. Fig.~1 shows the top-level block diagram. Data
flows from the fingerprint sensor inputs through the quality and
threshold pipeline, into the matching engine, and finally into the
authentication controller which makes the final access decision and
manages security responses.

The core design philosophy is \textit{hardware autonomy}: every
security action---quality checking, threshold adaptation,
template comparison, penalty enforcement, wipeout, and lockdown---
happens entirely in hardware logic. No processor, no operating system,
no software driver is involved in any security-critical path.

\subsection{Fingerprint Feature Representation}

Each fingerprint is represented as a compact 32-bit vector containing
four 8-bit fields: the X and Y coordinates of the fingerprint's core
point on the sensor surface, the dominant ridge orientation angle
(Theta), and the ridge density (how many ridges appear per unit area).
This four-dimensional representation, written as
$\{X[7:0], Y[7:0], \Theta[7:0], \rho[7:0]\}$,
captures the spatial and structural characteristics of a fingerprint
in a format that hardware arithmetic can process in a single clock
cycle. Feature vectors are extracted from the SOCOFing fingerprint
dataset using an OpenCV-based Python tool that processes real sensor
images and outputs 32-bit hex values ready to load into the vault.

\subsection{Quality Analyzer Module}

Before any matching is attempted, the system evaluates the quality
of the incoming fingerprint scan. Three sensor signals are read:
the surface noise level (how much interference is on the sensor,
caused by sweat, dust, or oil), the ridge contrast (how clearly
defined the fingerprint ridges are), and the spatial uniformity
(how consistently the ridge pattern flows across the sensor area).

These three values are combined into a single 8-bit quality score
between 0 (very poor) and 255 (excellent). The quality score is
then classified into one of four grades: PRISTINE (clean, dry,
high-contrast scan), GOOD (minor noise, still reliable), DEGRADED
(wet, dry, or aged skin with reduced clarity), or REJECT (score
too low, or contrast suspiciously flat as seen in artificial finger
spoofs). If the scan is classified as REJECT, the system immediately
denies access without running the matcher at all.

The anti-spoof guard floor is embedded within the quality grading:
if ridge contrast falls below a minimum value of 30 (out of 255),
the scan is unconditionally rejected regardless of any other quality
metric. Real human fingers always produce some contrast; flat
gelatin or silicone artificial fingers produce contrast readings
near zero.

\subsection{Dynamic Threshold Calculator Module}

Traditional biometric matchers compare the probe-to-template distance
against a single fixed number. If the distance is below that number,
access is granted; if above, it is denied. The problem is that this
fixed number cannot simultaneously serve all scan conditions: a
threshold tight enough to block impostors will also reject genuine
users with degraded-quality scans.

Our threshold calculator solves this by computing the acceptance
threshold dynamically from the quality score. When quality is high
(clean scan), the threshold is set small---meaning the probe must
closely match the template. When quality is lower (wet or dry finger),
the threshold is relaxed---the probe is allowed a larger distance from
the template before being rejected. This relationship is computed
using a simple bit-shift operation that requires no dividers or
multipliers, consuming zero DSP blocks on the FPGA. The formula is:

\begin{equation}
T(Q) = \min\bigl(28,\;\; 8 + \lfloor(255 - Q) / 16\rfloor\bigr)
\end{equation}

A quality score of 245 (pristine) gives a threshold of 8.
A quality score of 80 (wet finger) gives a threshold of 18.
Spoofed or rejected scans force the threshold to zero,
guaranteeing rejection. Table~\ref{tab:threshold} lists key
operating points.

\begin{table}[htbp]
\caption{Threshold Values at Key Quality Levels}
\begin{center}
\begin{tabular}{|l|c|c|c|}
\hline
\textbf{Scan Condition} & \textbf{Quality} & \textbf{Threshold} & \textbf{Grade} \\
\hline
Perfect, dry, clean      & 245 & 8  & PRISTINE \\
Normal everyday scan     & 200 & 11 & PRISTINE \\
Slightly moist           & 155 & 14 & GOOD     \\
Wet or sweaty finger     & 80  & 18 & DEGRADED \\
Very degraded            & 50  & 22 & DEGRADED \\
Artificial finger / spoof & --- & 0  & REJECT   \\
\hline
\end{tabular}
\label{tab:threshold}
\end{center}
\end{table}

\subsection{Secure Template Vault Module}

The secure template vault stores the enrolled fingerprint vectors for
up to four registered users. During normal operation it works like a
read-only memory: given a user ID, it outputs that user's stored
feature vector to the matcher. The enrolled data is loaded from
real fingerprint images processed through the SOCOFing dataset.

The vault's key innovation is its self-destruction capability. When
the authentication controller raises the \texttt{zeroize\_trigger}
signal, the vault immediately executes a three-pass memory purge:

\textbf{Pass 1 --- All Ones:} Every bit in every stored template is
set to 1 (\texttt{0xFFFFFFFF}). This saturates the storage cells,
driving them to a known maximum state.

\textbf{Pass 2 --- All Zeros:} Every bit is then set to 0
(\texttt{0x00000000}). This discharges all storage nodes. The
alternating write pattern is specifically required by NIST SP 800-88
and DoD 5220.22-M media sanitization guidelines to prevent residual
data remanence.

\textbf{Pass 3 --- Random Noise:} All storage locations are then
overwritten with pseudo-random bit patterns generated by an on-chip
16-bit Galois Linear Feedback Shift Register (LFSR). The LFSR
is initialized with a fixed seed and runs continuously; its output
is different for each location, preventing any pattern from remaining
in the memory cells.

After all three passes complete, the vault signals that templates are
permanently wiped. The chip then enters LOCKDOWN mode and never exits,
regardless of what signals arrive afterward. The vault also monitors
an inactivity timer: if the TTL (Time-to-Live) counter expires without
a successful authentication, the same three-pass purge runs
automatically---protecting against long-term physical access attacks
even if no active login attempt is ever made.

\subsection{Matcher Engine Module}

The matcher engine computes how different the incoming probe
fingerprint is from the enrolled template. Rather than comparing
individual feature fields one at a time, all four dimensions are
compared simultaneously and their differences are summed into a
single distance value. This approach---called 4D Manhattan
distance---adds the absolute differences in X position, Y position,
ridge angle, and ridge density:

\begin{equation}
D = |\Delta X| + |\Delta Y| + |\Delta\Theta| + |\Delta\rho|
\end{equation}

A small value of D means the probe closely matches the stored
template---likely the genuine user. A large value means the probe
is very different---likely an imposter. The matcher then compares
D against the dynamically computed threshold T(Q): if $D \leq T(Q)$,
access is granted; otherwise it is denied. All of this computation
happens in a single clock cycle using purely combinational logic---
no registers, no pipeline delays, no clock cycles wasted waiting.
The maximum possible distance value is $4 \times 255 = 1020$,
which the 10-bit accumulator handles easily.

\subsection{Authentication Controller FSM}

The authentication controller is the central brain of the system.
It is implemented as a Moore Finite State Machine (FSM) with seven
states, each representing a distinct phase of the authentication
process. Fig.~2 illustrates the FSM.

\textbf{S\_IDLE:} The system is waiting for someone to start a
scan. As soon as the start signal arrives, it moves to FETCH.
If the tamper pin is asserted at any time in any state, the FSM
immediately jumps to ZEROIZE.

\textbf{S\_FETCH:} The controller reads the enrolled template from
the vault for the requested user ID and presents it to the matcher.

\textbf{S\_COMPARE:} The matcher computes the 4D distance. The
quality analyzer simultaneously computes the quality score and
the threshold calculator computes T(Q). All three results are
ready in the same clock cycle.

\textbf{S\_EVALUATE:} The controller checks whether the distance
is within the threshold. If yes, access is granted and the fail
counter is reset. If no, access is denied, the fail counter
increments, and the system moves to PENALTY (if fails remain)
or directly to ZEROIZE (if this was the third failure).

\textbf{S\_PENALTY:} A configurable cooldown timer prevents rapid
repeated attempts. The system waits here for the penalty duration
before returning to IDLE for the next attempt. This makes
brute-force attacks extremely slow even if the attacker somehow
avoids triggering the three-strike limit.

\textbf{S\_ZEROIZE:} The controller fires the zeroize trigger,
which causes the vault to execute its three-pass purge. The
tamper alert output is asserted to notify any external monitoring
system. After zeroization completes, the system moves permanently
to LOCKDOWN.

\textbf{S\_LOCKDOWN:} Terminal state. The system self-loops here
forever. No authentication attempt, reset signal, or power cycle
(without full hardware reprogramming and re-enrollment) can exit
this state.

\subsection{Top-Level Integration (biometric\_top.v)}

The top-level module connects all five sub-modules with the
appropriate signal wires. It exposes all inputs---the scanned
vector, three quality signals, control inputs, user ID, and clock---
and all outputs including access granted/denied flags, fail count,
quality grade, distance value, tamper alert, templates wiped flag,
and lockdown status. All signals are accessible for waveform
monitoring during simulation.

% ============================================================
%  SECTION IV: SIMULATION AND VERIFICATION
% ============================================================
\section{Simulation and Verification}

The design was simulated using Icarus Verilog 0.10 with waveform
output captured to a VCD file, viewable in GTKWave. An equivalent
simulation was also run on EDA Playground using the same tool.

The testbench \texttt{tb\_biometric\_top.v} automatically runs
seven distinct scenarios and checks results without manual inspection.
Table~\ref{tab:sim} summarizes all seven tests. User~0's enrolled
template is $\{X=120, Y=145, \Theta=88, \rho=64\}$.

\begin{table}[htbp]
\caption{Testbench Scenarios and Results}
\begin{center}
\begin{tabular}{|c|p{2.9cm}|c|c|c|}
\hline
\textbf{\#} & \textbf{Scenario} & \textbf{Quality} & \textbf{T(Q)} & \textbf{Result} \\
\hline
1 & Clean scan, User 0 ($D\!=\!3$)          & 210 & 10 & GRANTED \\
\hline
2 & Wet finger, User 0 ($D\!=\!16$)         & 80  & 18 & GRANTED$^{\dagger}$ \\
\hline
3 & Flat spoof (contrast$=\!15$)            & --- & 0  & BLOCKED$^{\ddagger}$ \\
\hline
4 & Imposter 1 ($D\!=\!149$)                & 235 & 8  & DENIED \\
\hline
5 & Imposter 2 ($D\!=\!128$)                & 235 & 8  & DENIED \\
\hline
6 & Imposter 3 --- Strike 3                 & 235 & 8  & ZEROIZE \\
\hline
7 & Genuine user after lockdown             & --- & ---  & BLOCKED \\
\hline
\multicolumn{5}{l}{$^{\dagger}$A fixed threshold of 8 would have falsely rejected this user.}\\
\multicolumn{5}{l}{$^{\ddagger}$Anti-spoof guard triggers before matcher runs.}\\
\end{tabular}
\label{tab:sim}
\end{center}
\end{table}

\textbf{Test 1 (Pristine scan):} The probe is nearly identical to
the enrolled template (distance = 3). Quality is high (210), so the
threshold is set to 10. Since 3 $\leq$ 10, access is granted
immediately. This confirms basic matching works correctly.

\textbf{Test 2 (Wet finger --- key novelty demonstration):} The user
presents a sweaty finger. The probe is still the same genuine user but
the scan is noisier, resulting in a larger computed distance (16) and
a lower quality score (80). With a fixed threshold of 8, this user
would be rejected---a false rejection of a legitimate person. With our
dynamic threshold, quality score 80 produces T(Q) = 18, and since
16 $\leq$ 18, the genuine user is correctly accepted. This is the
central demonstration of Contribution 2.

\textbf{Test 3 (Artificial spoof finger):} A gelatin or silicone
artificial finger is presented. Its ridge contrast reading is 15,
which is below the anti-spoof floor of 30. The quality analyzer
immediately forces quality\_acceptable low. The matcher never runs.
Access is denied before any comparison occurs. This demonstrates
Contribution 3.

\textbf{Tests 4--5 (Imposter attempts):} Two imposter probes are
submitted. Both have large distances (149 and 128 respectively) from
the enrolled template---far above the threshold of 8. Both are denied.
The fail counter increments to 1, then 2.

\textbf{Test 6 (Strike 3 --- zeroization):} The third imposter attempt
fails. The fail counter reaches 3 (the maximum). The FSM transitions
to S\_ZEROIZE. The vault executes all three overwrite passes.
\texttt{templates\_wiped} is asserted, \texttt{lockdown} is asserted.
This is the central demonstration of Contribution 1.

\textbf{Test 7 (Post-lockdown genuine scan):} The real User~0 now
attempts to log in with a perfect scan---identical to the stored
template, distance = 0. Since the chip is in permanent lockdown and
the templates have been wiped, access is still denied. This confirms
the lockdown is irreversible and the zeroization was complete.

% ============================================================
%  SECTION V: FPGA SYNTHESIS RESULTS
% ============================================================
\section{FPGA Synthesis Results}

The complete SoC design was synthesized using Xilinx Vivado ML Edition
2024.1 targeting the AMD Artix-7 FPGA (part number xc7a35tcsg324-1,
speed grade $-1$), which is the device present on the Digilent Basys~3
development board. The module \texttt{biometric\_top} was set as the
top-level synthesis entity. Default Vivado synthesis and implementation
strategies were used with no manual timing constraints.

\begin{table}[htbp]
\caption{FPGA Synthesis Results (AMD Artix-7 xc7a35tcsg324-1)}
\begin{center}
\begin{tabular}{|l|c|r|}
\hline
\textbf{Resource / Metric} & \textbf{Value} & \textbf{Utilization} \\
\hline
Slice LUTs (Look-Up Tables) & 82   & 0.40\% of 20,800 \\
Slice Registers (Flip-Flops)& 51   & 0.12\% of 41,600 \\
DSP Blocks                  & 0    & 0\% of 90        \\
Block RAM Tiles             & 0    & 0\% of 50        \\
Max.\ Operating Frequency   & $\approx$175~MHz & ---  \\
Total On-Chip Power         & 62~mW & ---              \\
\quad\quad Dynamic Power    & 41~mW & ---              \\
\quad\quad Static (Leakage) & 21~mW & ---              \\
\hline
\end{tabular}
\label{tab:synth}
\end{center}
\end{table}

The most significant result in Table~\ref{tab:synth} is the zero
DSP block and zero block RAM utilization. DSP blocks are specialized
multiplier-accumulator units; using zero of them confirms that the
dynamic threshold calculation (which could have required division)
was successfully implemented using only bit-shift operations in
general-purpose LUT logic. The three-pass zeroization loop, including
the LFSR noise generator, also fits entirely within the FSM register
file without needing any dedicated memory.

The low LUT count (82 out of 20,800 available, or 0.4\%) means this
entire biometric authentication SoC occupies less than half a percent
of the smallest Artix-7 device. It could realistically be embedded
alongside other logic in a much larger design---for example, in a
complete IoT security subsystem or smart card controller.

The maximum operating frequency of approximately 175~MHz is derived
from the worst-case timing slack reported by Vivado. This is
calculated as $f_{max} = 1 / (T_{clk} - WNS)$ where $T_{clk}$ is
the target clock period (10~ns for 100~MHz) and WNS is the worst
negative slack margin. The critical timing path runs through the
4D Manhattan distance accumulator in the matcher engine.

The reference design in~\cite{verilog_biometric} reported a total
on-chip power of 0.184~W for its biometric DES co-processor on a
similar Xilinx FPGA. Our design consumes 62~mW (0.062~W), a
reduction of approximately 66\%, while providing significantly
stronger security guarantees through autonomous zeroization and
adaptive threshold matching.

% ============================================================
%  SECTION VI: DISCUSSION
% ============================================================
\section{Discussion}

\subsection{Why Hardware, Not Software?}

The two novel features in this paper---self-destructing templates
and quality-adaptive thresholds---could in principle be implemented
in software running on a processor. However, software implementations
have critical weaknesses in a security context. A software-based
zeroization routine can be bypassed, suspended, or subverted by an
attacker who has physical access to the system and can manipulate
power, clock, or memory signals. Fault injection attacks on
cryptographic software have been well documented~\cite{fault_attack};
the same attacks apply to software-based secure deletion.

A hardware FSM-based zeroization has no software stack to attack.
Once the \texttt{zeroize\_trigger} is raised, the purge sequence runs
at the hardware clock rate---in this design, completing in under
20~ns per pass---with no way for software to interrupt it. Similarly,
the quality-adaptive threshold computed in combinational logic is
instantaneous and cannot be tampered with through software means.

\subsection{Comparison with PUF-Based Security}

PUF-based authentication systems~\cite{puf_auth} generate a device
fingerprint from manufacturing variation, enabling zero-storage
device authentication. Our approach is complementary: where PUFs
protect device identity, our vault protects user biometric identity.
A future extension of this work could combine both---using a PUF to
generate the LFSR seed for the random noise pass, making the noise
pattern in Pass~3 of the zeroization unique to each chip instance and
thus resistant to precomputed pattern removal.

\subsection{Limitations and Future Work}

The current design stores templates for four users and matches a
single probe against one user at a time. Scaling to larger user
populations would require a wider address space and a sequential
search or hashing scheme. The feature vector is 32 bits representing
a single minutia point; real fingerprint matching typically uses
sets of 10--100 minutiae. Future work will extend the vault to store
multi-minutia sets and implement a hardware pipeline for set-level
matching. Additionally, as suggested by PUF literature~\cite{puf_auth},
integrating a ring oscillator PUF as a hardware root of trust for the
initial template enrollment process would further strengthen the
system against enrollment-time attacks.

% ============================================================
%  SECTION VII: CONCLUSION
% ============================================================
\section{Conclusion}

This paper presented a complete Verilog biometric authentication SoC
addressing two security problems that no prior hardware implementation
has solved: permanent on-chip template exposure and fixed matching
thresholds that fail under varying finger conditions. The
self-destructing template vault autonomously executes a three-pass
NIST-compliant memory wipeout and enters permanent hardware lockdown
upon detecting three failed login attempts, a tamper signal, or an
inactivity timeout---without any software involvement. The quality-adaptive
dynamic threshold engine reads three real-time sensor quality signals
and continuously adjusts the matching acceptance boundary, eliminating
wet-finger false rejections while maintaining strict imposter rejection.
An anti-spoof guard unconditionally blocks artificial finger attacks
by checking ridge contrast before the matcher runs. Synthesized on
AMD Artix-7 FPGA, the design achieves 175~MHz at 62~mW using only
82 LUTs and 51 flip-flops with zero DSP or BRAM consumption---a 66\%
power reduction compared to prior Verilog biometric hardware. To the
authors' knowledge, both the self-destructing vault and the
hardware quality-adaptive threshold are presented here for the first
time in any published biometric hardware design.

% ============================================================
%  ACKNOWLEDGMENT
% ============================================================
\section*{Acknowledgment}

The authors thank the Department of Electronics and Computer Science
Engineering, Vidyalankar Institute of Technology, Mumbai, for
providing FPGA hardware resources and EDA software licenses used
in this work.

% ============================================================
%  REFERENCES
% ============================================================
\begin{thebibliography}{00}

\bibitem{verilog_biometric}
Author(s),
``Verilog implementation of biometric based transmission,''
\textit{[Journal/Conference name]},
[Year].
\textit{[Replace with exact citation from the PDF in your folder]}

\bibitem{puf_auth}
Author(s),
``Physical unclonable functions for device authentication,''
\textit{[Journal/Conference name]},
[Year].
\textit{[Replace with exact citation from the PDF in your folder]}

\bibitem{secret_key_ic}
Author(s),
``Secret key generation in integrated circuits,''
\textit{[Journal/Conference name]},
[Year].
\textit{[Replace with exact citation from the PDF in your folder]}

\bibitem{fault_attack}
Author(s),
``Fault attack on supersingular cryptosystems,''
\textit{[Journal/Conference name]},
[Year].
\textit{[Replace with exact citation from the PDF in your folder]}

\bibitem{jain2004}
A.~K. Jain, A.~Ross, and S.~Prabhakar,
``An introduction to biometric recognition,''
\textit{IEEE Trans. Circuits Syst. Video Technol.},
vol.~14, no.~1, pp.~4--20, Jan. 2004.

\bibitem{nistsp80088}
M.~S. Kissel \textit{et al.},
``Guidelines for media sanitization,''
NIST Special Publication 800-88 Rev.\ 1,
National Institute of Standards and Technology,
Gaithersburg, MD, USA, Dec. 2014.

\bibitem{jiang2007}
R.~Jiang and D.~Crookes,
``An FPGA-based parallel matching processor architecture for fingerprint
minutiae,''
in \textit{Proc. IEEE Int. Symp. Circuits Syst. (ISCAS)},
pp.~3314--3317, 2007.

\bibitem{ratha2001}
N.~K. Ratha, J.~H. Connell, and R.~M. Bolle,
``Enhancing security and privacy in biometrics-based authentication
systems,''
\textit{IBM Syst. J.}, vol.~40, no.~3, pp.~614--634, 2001.

\end{thebibliography}

\end{document}
